% Replace the existing Section 3.3 with this subsection.
% This file is a manuscript fragment, not a standalone compilable document.
\subsection{Closed-loop rollout training}
\label{sec:rollout_training}

Each specialist is initialized by encoding a physical history into
its own local coordinates. It then advances autonomously, using its
predicted velocity and reconstructed pressure to update the history.
No reference state inside the prediction window is injected after
initialization.

\paragraph{Frozen-context advancement.}
During one reduced velocity step, the pressure $\bm b_n^{(r)}$,
preceding-state history, and parameter define a fixed context
$\mathcal C_n^{(r)}$. For a trial velocity state $\bm a$, the hybrid
right-hand side and numerical update are
\begin{align}
  F_r(\bm a;\mathcal C_n^{(r)})
  &=\mathcal F_r^G(\bm a,\bm b_n^{(r)};\mu)
    +\bm\rho^u_{\theta_r}\!\left(
      \bm\xi^{(r)}(\bm a;\mathcal C_n^{(r)})\right),
    \label{eq:main_frozen_context_rhs}
    \\
  \bm a_{n+1}^{(r)}
  &=\operatorname{Step}_r\!\left(
     \bm a_n^{(r)};
     F_r(\cdot;\mathcal C_n^{(r)}),\Delta t_r\right).
    \label{eq:main_specialist_step}
\end{align}
Here $\Delta t_r$ is the reduced time step and
$\operatorname{Step}_r$ denotes the chart-specific advancement rule.
For continuous-time instances, Runge--Kutta stages evaluate the
velocity-dependent right-hand side with the same frozen context;
the feature vector is evaluated at the trial velocity, rather than
held entirely fixed. The updated velocity is then supplied to
Eq.~\eqref{eq:main_pressure_closure}, and the local history is shifted
before the next step. The nominal Runge--Kutta order applies to the
frozen-context velocity subproblem, not automatically to the complete
pressure/history-coupled scheme. Benchmark-specific stepping rules,
including the data-only exception, are given in
Appendix~\ref{app:implementation}.

\paragraph{Rollout objective.}
Training unrolls the complete specialist update for $K$ steps and
minimizes the normalized coefficient error
\begin{equation}
  \mathcal L_r=\frac{1}{K}\sum_{k=1}^{K}\left[
  \frac{\left\|(\widehat{\bm a}_{n+k}^{(r)}-\bm a_{n+k}^{(r)})
  \oslash\bm\sigma_a^{(r)}\right\|_2^2}{d_u^{(r)}}
  +\frac{\left\|(\widehat{\bm b}_{n+k}^{(r)}-\bm b_{n+k}^{(r)})
  \oslash\bm\sigma_b^{(r)}\right\|_2^2}{d_p^{(r)}}\right].
  \label{eq:main_rollout_loss}
\end{equation}
Hats denote recursively predicted coefficients, while unhatted
coefficients are reference targets. The dimensions $d_u^{(r)}$ and
$d_p^{(r)}$ are chart-specific; the scaling vectors
$\bm\sigma_a^{(r)}$ and $\bm\sigma_b^{(r)}$ are estimated from the
corresponding training set, and $\oslash$ denotes componentwise
division. Reference coefficients are used to compute the loss but
are not fed back into the predicted trajectory. Progressively longer
rollout horizons expose the learned corrections to their own
accumulated prediction errors. The curricula and optimization
settings are reported in Appendix~\ref{app:implementation}, and
held-out trajectory splits are listed in Table~\ref{tab:chart_splits}.

The training objective is distinct from physical-field evaluation:
variance normalization reweights reduced-coordinate errors and is
not a physical energy norm. Weighted POD orthonormality gives
$\|\Phi_u^{(r)}(\widehat{\bm a}-\bm a)\|_{M_u}^2
=\|\widehat{\bm a}-\bm a\|_2^2$ within a common chart, but comparison
with the original CFD field additionally includes its POD projection
remainder. The coefficient loss alone guarantees neither
conservation nor long-term stability.
